\documentclass[11pt,a4paper]{article}
\usepackage[utf8]{inputenc}
\usepackage[spanish]{babel}
\usepackage{times}
\usepackage{setspace}
\usepackage[margin=2.5cm]{geometry}

\onehalfspacing

\begin{document}

\begin{center}
    \textbf{\Large DOCUMENTO DE POSICIÓN OFICIAL}\\[0.5cm]
    \textbf{Modelo:} Modelo de las Naciones Unidas\\
    \textbf{Comité:} Consejo de Seguridad de la ONU (Histórico)\\
    \textbf{Tema:} La crisis en Rumania, eventos de Timișoara y cierre de fronteras\\
    \textbf{País Representado:} Reino de Noruega\\
    \textbf{Delegado/a:} [Nombre del Delegado]
\end{center}

\vspace{0.5cm}

\section*{Contexto de Política Exterior}
El Reino de Noruega, como miembro fundador de la Organización del Tratado del Atlántico Norte (OTAN) y firme defensor de los valores democráticos, ha mantenido una política exterior estructurada sobre el equilibrio entre la disuasión militar y la diplomacia pacificadora frente al bloque soviético. Durante la Guerra Fría, esta delegación adoptó una postura de confianza y reducción de tensiones para evitar la provocación directa hacia la Unión Soviética, priorizando la estabilidad subregional y la resolución pacífica de controversias. Hacia fines de 1989, el Reino de Noruega sostiene la firme convicción de que la soberanía nacional no puede ser esgrimida como escudo para la supresión sistemática de los derechos fundamentales. En concordancia con el Acta Final de Helsinki firmada en 1975, el Estado noruego entiende que la imperativa protección de los derechos humanos es una responsabilidad internacional compartida.

Esta delegación considera fundamental recalcar que, ante la reciente efervescencia democrática en Europa del Este, Noruega ha promovido enérgicamente las transiciones pacíficas, rechazando el uso de la fuerza desproporcionada contra la sociedad civil. Siendo un país con vasta tradición en la facilitación del diálogo y la protección humanitaria, el gobierno de Noruega condena las acciones estatales orientadas a coartar libertades legítimas. Se sostiene oficial y categóricamente que la jurisdicción interna debe operar en cumplimiento armónico con los compromisos universales y aquellos tratados vinculantes a los cuales los Estados han ofrecido su ratificación, por lo que toda transgresión masiva a la población merece ser repudiada e intervenida diplomáticamente.

\section*{Análisis Jurídico y Fáctico de la Crisis}
Los drásticos acontecimientos que subyacen a las movilizaciones de Timișoara en diciembre de 1989 y la letal represión ejercida por el gobierno de Nicolae Ceaușescu preocupan profundamente al Estado noruego. La información verificada respecto a fuerzas armadas abriendo fuego indiscriminado contra la población civil desarmada, sumada a la determinación de cerrar arbitrariamente las fronteras rumanas, configura un escenario de inestabilidad crítica. Si bien el Artículo 2, numeral 7 de la Carta de las Naciones Unidas protege la jurisdicción interna e inviolabilidad de los Estados (Organización de las Naciones Unidas [ONU], 1945), esta delegación argumenta jurídicamente que la grave supresión de derechos inalienables y el inminente éxodo de refugiados representan una amenaza a la paz y la seguridad internacional.

Por consiguiente, este asunto trasciende el ámbito estrictamente soberano, ingresando en las competencias del Consejo de Seguridad para la preservación de la paz y seguridad globales. Al vulnerar el gobierno rumano los principios fundamentales de los Pactos Internacionales de Derechos Civiles y Políticos de 1966, de los cuales es signatario, la evidencia fáctica demuestra que la violencia sistémica empleada contra su propio pueblo exige el escrutinio de este órgano. El Reino de Noruega concluye que ignorar la gravedad del aislamiento impuesto, justificándolo bajo el falso amparo de la no injerencia, es inaceptable y requiere la atención prioritaria de la comunidad internacional.

\section*{Propuestas de Resolución}
Para hacer frente a esta escalada represiva en Rumania, el Estado noruego impulsará dos líneas estratégicas ante el Consejo de Seguridad. Principalmente, se exigirá el cese definitivo de la hostilidad estatal y la inmediata reapertura de las fronteras nacionales para facilitar el ingreso de ayuda de la Cruz Roja, de manera expedita. Simultáneamente, se demandará la conformación de una misión civil neutral con observadores de la ONU o embajadores del entorno OSCE/CSCE, para verificar las violaciones en terreno y reportar formalmente a la asamblea. Como medida de presión complementaria, Noruega propondrá la congelación temporal de transferencias financieras no esenciales dirigidas al aparato gubernamental de Bucarest, sancionando diplomáticamente al régimen sin comprometer los suministros de primera necesidad de los civiles.

Para consolidar estas iniciativas, esta delegación fomentará alianzas proactivas con repúblicas de Europa Occidental, naciones nórdicas y miembros del bloque atlántico, fortaleciendo una respuesta multilateral sólida en defensa de la justicia civil. No obstante, el Reino de Noruega advierte formalmente sobre su línea roja inquebrantable: esta delegación se opondrá firmemente y jamás consentirá la implementación de intervenciones armadas foráneas de carácter militar sobre el territorio de Rumania. Dichas medidas contravienen explícitamente el uso proscrito de la fuerza estipulado en el Artículo 2, numeral 4 de la Carta de la ONU, optando en todo momento por forzar reformas mediante las vías de la pacificación y la máxima presión diplomática multilateral.

\vspace{1cm}
\begin{center}
    \textbf{Referencias}
\end{center}
\begin{description}
    \item Organización de las Naciones Unidas. (1945). \textit{Carta de las Naciones Unidas}. San Francisco.
    \item Organización de las Naciones Unidas, Comisión de Derechos Humanos. (1989). \textit{Resolución 1989/75. Situación de los derechos humanos en Rumania}. Ginebra.
\end{description}

\end{document}
